\scanpage{60}
condition, rappelons-le, que l'interprétation
d'un polygône comme une structure
\uline{sur $S$} (ens des sommets) -- savoir
un ens \unsure{de des|de de|des} 2 \unsure{paires|pourv.} \unsure{circulaire -- est|circulairement|circulaires -- est}
valable. \unsure{Par|Pour} la suite, on n'aura
à regarder que les cas $n\geq3$.

6. \uline{Orientation} d'un contour

a) Soit $R^+\subset R$, \deleted{\unsure{on a|on}} $R^-=R\smallsetminus R^+$. Conditions
équivalentes
\[
\begin{aligned}
\text{a)}\quad&\sigma_0R^+=\sigma_1R^+=R^-\\
\text{b)}\quad&R^+\xrightarrow{\sim} S,\quad R^+\simeq A
\end{aligned}
\quad\left|\begin{gathered}\text{On dit alors que }R^+\\\text{est une orientation}\end{gathered}\right.
\]
c'est ainsi une partie \inserted{$R^+$} de $R\deleted{\unsure{{}^+}}$ qui est le graphe
d'une \uline{bijection} $S\overset{\sim}{\longleftrightarrow} A$.
\margin{Orientation \deleted{\unsure{positive|opposée}} \inserted{opposée} d'une orientation \deleted{\unsure{déf}} (involution -- \unsure{ss|est} $\unsure{\ne\mathrm{id}|\ne 0}$ si $\Pi\ne\varnothing$)}
\[
\text{b)}\qquad\mathrm{Or}\left(\coprod_{i\in I}\deleted{\unsure{R_i(C)|R_i(C_i)}}\underset{\deleted{\text{contour}}}{C_i}\right)\simeq\prod_i\mathrm{Or}(C_i)
\]
si $(C_i)_{i\in I}$ est une famille de contours.

c) Si $\Pi$ est un polygône combinatoire,
alors $\mathrm{Card}\,\mathrm{Or}(\Pi)=2$, et $R$ est \inserted{alors} somme
disjointe de la \deleted{\unsure{ens|en}}famille des deux orientations
de $\Pi$, \deleted{\unsure{Pour tte orientation}} \inserted{lesquelles sont opposées l'une de l'autre.} Les deux \unsure{orient\textsuperscript{ns}|orientations}
sont les deux orbites de $\unsure{G_1^{0}|G_1^\circ|\mathfrak G_1^\circ}$ dans $R$.
% archive sheet number: 59
